\section{Problem setting and notation}\label{sec:setting}

\paragraph{Users, items and outcomes.}
There are $m$ users $u\in[m]$ and $n$ items $i\in[n]$. Each pair $(u,i)$ has a binary
potential outcome $Y_{ui}\in\{0,1\}$, which records whether user $u$ would like item
$i$ if it were shown. Each pair also has an exposure indicator $O_{ui}\in\{0,1\}$,
which records whether the pair was exposed. A pair is a logged positive when
$O_{ui}Y_{ui}=1$. The training data record the exposure indicator $O$ and the logged
outcome $R_{ui}=O_{ui}Y_{ui}$. The outcome $Y$ is not recorded on unexposed pairs. The
outcome $Y$ is what the recommender should estimate, and the exposure $O$ is the
mechanism that makes the log biased.

What $O$ records depends on the platform. On KuaiRec it is an impression chosen by the
system. On Coat and Yahoo!\,R3 it is a rating that the user chose to give, so the
``exposure'' modelled by the propensity is self-selection by the user rather than the
display policy of a platform. All statements below are relative to the mechanism that
generates $O$, and we keep the word exposure for brevity.

\paragraph{Data.}
The training data are a biased log $\mathcal D_{\mathrm{log}}=\{(u,i,O_{ui},O_{ui}Y_{ui})\}$
generated by the exposure mechanism. A separate uniformly collected sample
$\mathcal D_{\mathrm{unif}}$ contains outcomes for randomly exposed pairs. The uniform
sample is used only to select hyper-parameters and to evaluate models. It never enters
the training objective, and the propensity and imputation models are fitted on the log
alone. The one exception is the Coat dataset, whose shipped propensities were estimated
by its authors from a uniform sample \citep{schnabel2016}; Section~\ref{sec:discussion} returns to this point.

For each user $u$, the set of non-logged pairs is
\[
\mathcal N_u=\{\,i\in[n] : O_{ui}Y_{ui}=0\,\}.
\]
This set contains exposed items that the user did not choose and unexposed items. A
pairwise sampler draws negatives for a logged positive $(u,i)$ from the items of the
catalogue and redraws once when the sampled item is a logged positive. The negatives
are therefore approximately uniform on $\mathcal N_u$.

\paragraph{Exposure model.}
Let $x_{ui}$ denote the covariates of the pair that the nuisance models use, for example
user and item exposure rates. The exposure propensity is
$p_{ui}=\Pr(O_{ui}=1\mid x_{ui})$, and the imputation is
$\hat Y_{ui}\in[0,1]$, an estimate of $\Pr(Y_{ui}=1\mid x_{ui})$ fitted on the log.
Both are functions of the covariates only. The method relies on the following
modelling assumption.
\begin{assumption}[Exposure independent of the outcome given covariates]\label{as:indep}
Conditionally on $x_{ui}$, the exposure $O_{ui}$ and the outcome $Y_{ui}$ are independent.
\end{assumption}
Under this assumption the joint distribution of $(O_{ui},Y_{ui})$ factorises given
$x_{ui}$, and the posterior $Q_{ui}$ of Eq.~\eqref{eq:Q} follows from Bayes' rule. The
quantity $Q_{ui}=\Pr(Y_{ui}=1\mid R_{ui}=0,\,x_{ui})$ is the probability that a
non-logged pair, which is the only kind of pair sampled as a negative, is a positive. It
is conditional on the covariates and on $R_{ui}=0$ only. It does not condition on $O_{ui}$.
Assumption~\ref{as:indep} excludes exposure that depends on $Y$ beyond the covariates. It
is the condition under which the nuisance models of this paper can be interpreted as
posterior quantities, and it is discussed as a limitation in Section~\ref{sec:discussion}.

To keep the inverse weights finite, the propensity is clipped from below at a floor
$\tau>0$, and we write $\bar p_{ui}=\max(\hat p_{ui},\tau)$ for the clipped estimate. The
experiments use $\tau=0.05$.

\paragraph{Evaluation target.}
At test time we rank the items of each user. For user $u$ the candidates are the items
that the user never interacted with in the log, and the relevance of a candidate is its
outcome $Y_{ui}$ as observed in the uniform test sample. The accuracy measure is a ranking metric,
nDCG$@K$, computed on this uniform sample. The cut-off $K$ depends on the dataset and
is given in Section~\ref{sec:setup}. Because these candidates have no logged positive, the
evaluation asks about preferences that the log does not reveal directly, which is the
quantity that the debiasing methods are meant to recover. Tests and model selection are
also described in Section~\ref{sec:setup}.

\paragraph{Scores.}
A recommender produces a score $s_{ui}$ for each pair. For a graph model with parameters
$\theta$, we write $f_\theta(u,i)$ for its propagated score. The offset added to this
score and the weights of the negatives are defined in Section~\ref{sec:method}. The edge
estimates used by the DRUP baseline are defined there as well, since only that baseline
needs them.
